\paragraph{7. Proof of Theorem 4.}
We can assume that $q\leq\log m$, for else the statistical statistical bound is trivial and the theorem is true. We apply the separator Theorem 14 with $\alpha=1/1000$ and the sequence
\[t_i:=m/w^{c_0(q/\alpha)^{-c_1i}},\]
for constants $c_0,c_1$ to be set later. For large enough $c_1$ this satisfies the hypothesis of the theorem that $t_i\geq t_{i-1}\cdot cqw/\alpha$. The hypothesis that $1-2^{-t_0}\geq\sqrt{8/|S|}=\sqrt{8/2^m}$ holds as well since $|S|=2^m$.
Let $k$ and $D_k$ be as given by the theorem. Let
\[X:=\operatorname{Rank}^{-1}f(D_k).\]
Note that this is a valid definition because $f(D_k)$ is in the range of $\operatorname{Rank}$, and the latter is 1-1. The separator theorem guarantees that $H_\infty(f(D_k))\geq m-t'_{k-1}$, where
\[t'_{k-1}:=t_{k-1}\cdot O(q^2w^2/\alpha^4).\]
Hence also $H(X)\geq m-t'_{k-1}$.
\paragraph{Comets.}
We now apply the comet-finding Lemma 21 to the $t_k=m/w^{c_0(q/\alpha)^{-c_1k}}$ trees given by the separator. For $c=1$, the lemma gives a set of
\[t'_k:=m/w^{c_0(q/\alpha)^{-c_1k}+O(1)}\]
disjoint $w$-comets. We shall only use that they are 100-comets, and their head lengths will not be relevant now. We want to find a comet whose outputs are ``sufficiently random.''
Define $a:=t'_{k-1}$ and $b:=t'_k$.
Partition $X$ into $b$ consecutive blocks, where each block contains exactly one comet and intersects no others. Let $Z_1,Z_2,\ldots,Z_b$ be the blocks, and let $|Z_i|=s_i$ with $\sum_i s_i=m$. Applying Lemma 22 we find $\geq b-a/\epsilon$ blocks $i$ such that $H(Z_i\mid Z_1Z_2\ldots Z_{i-1})\geq s_i-\epsilon$. We set $\epsilon=1/w$ (a sufficiently small constant would be enough), and we verify that $b-a/\epsilon\geq1$, yielding at least one block $i^*$ such that
\[H(Z_{i^*}\mid Z_1Z_2\ldots Z_{i^*-1})\geq s_{i^*}-\epsilon.\tag{1}\]
The inequality $b-a/\epsilon\geq1$ is true because we only lost $w^{O(1)}$ factors, so it suffices to make the constant $c_1$ large enough in the definition of the sequence $t_i$. Formally,
\[\frac ba=\frac{t_k}{t_{k-1}}\cdot\frac1{O(q^2w^2/\alpha^4)\cdot w^{O(1)}}\geq\frac{w^{c_1}}{w^{O(1)}}>w^{100}.\]
The inequalities holds for $c_1$ large enough and using $q\leq\log m$, $w\geq\log m$, $\alpha=\Theta(1)$.
\paragraph{Breaking correlation in the random comet.}
Hence we now have a comet $(p,i,j)$ that is contained in an interval $Z_{i^*}$ such that:
\begin{enumerate}
\item Equation 1 holds, and
\item $f_i(D_k),f_j(D_k)$ are $\alpha$-independent.
\end{enumerate}
The next lemma directly contradicts this and concludes the proof.
\paragraph{Lemma 23.}
Let $X_1,X_2,\ldots,X_m$ be $0$--$1$ random variables, and $(p,i,j)$ a $c$-comet for a sufficiently large $c$. Let $\ell:=i-p$ and $d:=j-i$.
Suppose that
\[H(X_{p+1},X_{p+2},\ldots,X_j\mid X_1,X_2,\ldots,X_p)\geq\ell+d-1/c.\]
Then there exists an integer $t$ such that
\begin{align*}
\mathbb P_X[\operatorname{Rank}(X)_j\geq t+\ell/2+d/2+c^{1/3}\sqrt d]&\geq1/10,\text{ and}\\
\mathbb P_X[\operatorname{Rank}(X)_i<t+\ell/2]&\geq1/10,\text{ but}
\end{align*}
\begin{multline*}
\mathbb P_X[\operatorname{Rank}(X)_j\geq t+\ell/2+d/2+c^{1/3}\sqrt d\ \wedge\ \operatorname{Rank}(X)_i<t+\ell/2]\\
\leq1/1000\quad(\ll1/10\cdot1/10).
\end{multline*}
\begin{proof}
Let us start with the last inequality, because we can prove it without getting our hands on $t$. The probability is at most
\[\mathbb P_X\left[\sum_{k=i+1}^jX_k\geq d/2+c^{1/3}\sqrt d\right].\]
By Pinsker's inequality ([14], Chapter 3; Exercise 17) the distribution of $X_{i+1},X_{i+2},\ldots,X_j$ is $4/\sqrt c$ close to the uniform $U_1U_2\ldots U_d$. Hence the above probability is
\[\leq\Pr_U\left[\sum_{k=1}^dU_k\geq d/2+c^{1/3}\sqrt d\right]+4/\sqrt c\leq1/2000+4/\sqrt c\leq1/1000.\]
where the second inequality follows from Chebyshev's inequality for sufficiently large $c$.
We now verify the first two inequalities in the conclusion of the lemma. Let $Y:=X_1,X_2,\ldots,X_p$ stand for the prefix, and $Z:=X_{p+1},X_{p+2},\ldots,X_j$ for the $\ell+d$ high-entropy variables. Let
\[A:=\{y\in\{0,1\}^p:H(Z\mid Y=y)\geq\ell+d-2/c\}\]
be the set of prefix values conditioned on which $Z$ has high entropy. We claim that $\mathbb P[Y\in A]\geq1/2$. This is because, applying Markov Inequality to the non-negative random variable $\ell+d-H(Z\mid Y=y)$ (for $y$ chosen according to $Y$),
\begin{align*}
\mathbb P[Y\notin A]&=\mathbb P_{y\in Y}[\ell+d-H(Z\mid Y=y)>2/c]\\
&\leq\mathbb E_{y\in Y}[\ell+d-H(Z\mid Y=y)]/(2/c)\\
&=(\ell+d-H(Z\mid Y))/(2/c)\leq(1/c)/(2/c)=1/2.
\end{align*}
Note that for every $y\in A$ we have, by definition, that the $(\ell+d)$-bit random variable $(Z\mid Y=y)$ has entropy at least $\ell+d-2/c$, and so by Pinsker's inequality ([14], Chapter 3; Exercise 17) the random variable $(Z\mid Y=y)$ is $(\epsilon:=4\sqrt{2/c})$-close to uniform over $\{0,1\}^{\ell+d}$. Therefore, for any subset $S\subseteq A$, the random variable
\[(Z\mid Y\in S)\text{ is }\epsilon\text{-close to uniform over }\{0,1\}^{\ell+d}.\tag{2}\]
Now define $t$ to be the largest integer such that
\[\mathbb P[Y\in A\wedge\operatorname{Rank}(X)_p\geq t]\geq1/4.\tag{3}\]
Since by definition of $t$ we have $\mathbb P[Y\in A\wedge\operatorname{Rank}(X)_p\geq t+1]<1/4$, we also have
\[\mathbb P[Y\in A\wedge\operatorname{Rank}(X)_p\leq t]\geq1/2-1/4=1/4.\tag{4}\]
We obtain the desired conclusions as follows, denoting by $U_1,U_2,\ldots,$ uniform and independent $0$--$1$ random variables. The first probability in the conclusion of the lemma is at least
\[\mathbb P[\operatorname{Rank}(X)_j\geq t+(\ell+d)/2+\sqrt\ell/c^{1/6}]\]
because $\ell\geq c\cdot d$.
Writing $\operatorname{Rank}(X)_j$ as the sum of the first $p$ bits and the rest, the above probability is at least
\begin{multline*}
\mathbb P\left[\sum_{k\leq\ell+d}Z_k\geq(\ell+d)/2+\sqrt\ell/c^{1/6}\mid Y\in A\wedge\operatorname{Rank}(X)_p\geq t\right]\\
\cdot\mathbb P[Y\in A\wedge\operatorname{Rank}(X)_p\geq t].
\end{multline*}
The second factor is $\geq1/4$ by (3). Also by (2) in the first factor we can replace the $Z_k$ with uniform bits changing the probability by at most $\epsilon$. Hence the first factor is at least
\[\mathbb P\left[\sum_{k\leq\ell+d}U_k\geq(\ell+d)/2+\sqrt\ell/c^{1/6}\right]-\epsilon.\]
In turn the probability is
\[\geq1/2-\sqrt\ell/c^{1/6}\cdot\Theta(1/\sqrt\ell)\geq1/2-\Theta(1/c^{1/6})\]
using an estimate of the central binomial coefficient provided e.g. in [13], Lemma 17.5.1.
Overall, the first probability in the conclusion of the lemma is
\[(1/2-\Theta(1/c^{1/6})-\epsilon)(1/4)\geq1/10\]
for large enough $c$.
We now turn to the second probability in the conclusion of the lemma. Proceeding in a similar way, this probability is at least
\begin{align*}
&\mathbb P\left[\sum_{k\leq\ell}Z_k<\ell/2\mid Y\in A\wedge\sum_kY_k\leq t\right]\cdot\mathbb P\left[Y\in A\wedge\sum_kY_k\leq t\right]\\
&\geq\left(\mathbb P\left[\sum_{k\leq\ell}U_k<\ell/2\right]-\epsilon\right)\cdot(1/4)\\
&\geq(1/2-\epsilon)\cdot(1/4)\geq1/10
\end{align*}
for all sufficiently large $c$. Here the second inequality uses (2) and (4).
\end{proof}
